\section{The shortening interface for bounded equivalence}\label{app:shortening}

This appendix gives the substantive sufficiency step of
Proposition~\ref{prop:bounded}. It is a candidate application of Sela's
graded shortening theorem~\cite{sela}; the JSJ and Rips machinery is
imported, not verified by the finite stabilizer program.

\subsection{The relative lemma and its normalization}
Let $A=F(a_1,\ldots,a_k)$, $k\ge2$, and let $u\ne1$ lie in no
proper free factor of $A$. The required assertion is that for fixed
target rank $r\ge2$ there is $K$ such that every injection
$f:A\to F_r$ admits $c\in F_r$ and $\theta\in\Aut(A)$ fixing
$u$ with
\begin{equation}\label{eq:relative-short}
 \max_i|cf(\theta(a_i))c^{-1}|\le K\clen{f(u)}.
\end{equation}

Adjoin a generator $z$, set $G=A*\langle z\rangle$ and
$P=\langle u,z\rangle$. There is no nontrivial free splitting with
$P$ in one factor. Indeed the intersection of that free factor with
$A$ is a free factor of $A$ containing $u$, hence all of $A$; it
also contains $z$. The same intersection argument shows that every
automorphism fixing $P$ pointwise preserves $A$ and restricts to
$\Aut(A/u)$.

Conjugate $f(u)$ to a cyclically reduced word, then minimize
$\max_i|f(\theta(a_i))|$ over $\Aut(A/u)$. The integer minimum
exists. Failure of~\eqref{eq:relative-short} would give minimizing
injections $f_m$ with
\[
 L_m=|f_m(u)|=\clen{f_m(u)},\qquad
 M_m=\max_i|f_m(a_i)|>2^m(1+L_m).
\]
Extend to injections $h_m:G\to F_r*\langle z\rangle$ fixing $z$.
Use the fixed generating set $\{a_1,\ldots,a_k,u,z\}$, with the
parameters explicit. The maps are shortest under the graded modular
subgroup, because it fixes $P$ and every such precomposition restricts
to the minimization already made.

In Sela's Definitions 10.1--10.3, the proposed reduction regards this
relative group as rigid or solid. If the graded JSJ is trivial it is
rigid. Otherwise minimizing an embedding in its graded modular orbit
gives a constant shortening sequence with quotient $G$, which dominates
every quotient and is the maximal shortening quotient; constant
sequences are allowed in the source's convention. The no-relative-free-product
condition follows from injectivity and the free-factor argument above.
The displayed $2^m$ inequality then gives a flexible graded sequence.
These are the hypotheses that must be checked in applying the deep
proper-quotient conclusion, rather than replacing it by ordinary
bounded-parameter shortening.

Here is the normalization check needed to identify its canonical
limiting quotient. In the target tree let $o$ be the identity,
$D_m(x)$ the maximum displacement by the chosen generators, and
$\mu_m=\min_xD_m(x)$ over vertices. The axes of $f_m(u)$ and $z$
meet precisely at $o$. One of them has projection of $x$ equal to
$o$, so $D_m(x)\ge2d(x,o)$. A generator attaining $M_m$ gives
$M_m\le2d(x,o)+D_m(x)$. For a minimizing vertex $x_m$,
\begin{equation}\label{eq:minimum-compare}
 M_m/2\le\mu_m\le M_m,\qquad d(x_m,o)/\mu_m\le1/2.
\end{equation}
Moreover, for some $s\in\{u,z\}$ along a subsequence,
$o\in[x_m,h_m(s)x_m]$. Thus $o$ and all its fixed orbit translates
belong to finite orbit hulls in the minimum-displacement normalization;
bounded distance alone is not being used to assume their survival.
After rescaling by $\mu_m$, $P$ fixes the limiting point of $o$.

Pass to a subsequence with $|f_m(a_i)|=M_m$ for one fixed $i$ and
$M_m/\mu_m\to\rho\in[1,2]$. The three points
$h_m(a_i)o,h_m(za_i)o,h_m(z^{-1}a_i)o$ have unscaled pairwise
distances $2M_m+1,2M_m+1,2M_m+2$, by free-product normal forms.
They give a nondegenerate tripod with all legs of length $\rho$ in
the same canonical limit. An element in the action kernel fixes this
tripod. Sela's elementary tripod lemma places it eventually in the
kernels of the $h_m$~\cite[Lemma 1.3(iii)]{sela}; injectivity makes
it trivial. Hence the limiting quotient is $G$ itself. The flexible
inequality survives passage to a subsequence. This contradicts the
properness of flexible graded quotients of rigid or solid groups
\cite[Lemma 10.4(ii)]{sela}, giving~\eqref{eq:relative-short}.

Shortness here concerns the original maps with their fixed parameter
images. Conjugating to the minimizing basepoint is used for the limit;
it is not asserted to preserve based graded shortness. The separate
\record{problems/F38/normalization-supplement.md}{normalization supplement}
records the complete comparison with the source convention.

\subsection{From a finite stabilizer orbit to the length bound}
Suppose first that $u$ fills $A$ as a free factor and the stabilizer
has a finite orbit on $[v]$ in $A$, represented by words of lengths
at most $B$. Applying~\eqref{eq:relative-short} to an injection $f$
and then to a representative of $\theta^{-1}[v]$ gives
\[
 \clen{f(v)}\le BK\clen{f(u)}.
\]
No value of $K$ is required by the eventual decision algorithm.

In ambient rank at least three choose a minimal free factor $A$
containing $u$. If $\rank A\ge2$, every complementary letter $b$
can be Nielsen-twisted as $b\mapsto ba_1$ or $b\mapsto ba_2$ for
distinct basis letters $a_1,a_2$ of $A$. Finiteness of $H_u[v]$
forces bounded lengths under both iterations. For a cyclic word whose
non-$e$ letters are $t_1,\ldots,t_s$ and intervening $e$-runs are
$e^{k_i}$, the exact formula for $d\mapsto de$ is
\begin{equation}\label{eq:nielsen-exact}
 \clen{\sigma_{d,e}^n(w)}
 =s+\sum_i|k_i+n(\mathbf1_{t_i=d}-\mathbf1_{t_{i+1}=d^{-1}})|.
\end{equation}
Inverse non-$e$ neighbors cannot newly cancel: their intervening run
was nonempty and its coefficient of $n$ is zero. Thus zero growth
requires each positive $b$ to pair with a negative $b$ across a pure
$a_j$ run. It cannot satisfy this for both $a_1,a_2$ unless no $b$
occurs. Hence $v$ is conjugate into $A$, and the preceding bound applies
to each restricted ambient automorphism.

If $A=\langle a\rangle$ and the ambient rank is at least three,
use the maps $b\mapsto bc$, $c\mapsto cb$, $b\mapsto ba$, and
$x_j\mapsto x_jb$ for the remaining basis letters. Formula~\eqref{eq:nielsen-exact}
shows that their common bounded conjugacy classes are exactly powers
of $a$. One explicit check is the product of the first three two-vertex
zero-growth graphs: its $b,c$ edges form a forest, leaving only
$a$-loops and loops of the unused generators after pruning. The last
maps remove those unused letters. The full eight-vertex graph is
listed in the retained \record{problems/F38/primitive-power-proof.md}{Nielsen argument}.
Therefore $v$ is a power of the same primitive element, with constant
length ratio.

Finally take ambient rank two and $u=a^p$. Bounded iterates under
$b\mapsto ba$ put $v$, up to conjugacy, in
$\langle a,bab^{-1}\rangle$. Fix an expression of length $L$ in
these generators. Under an automorphism put $A=\alpha(a)$,
$B=\alpha(b)$, $C=BAB^{-1}$ and $\ell=\clen A\ge1$.
The rank-two commutator identity gives
$\clen{AC^{-1}}=\clen{ABA^{-1}B^{-1}}=4$.
If the axes of $A,C$ are disjoint their distance is at most one,
by the tree product formula. At a nearest vertex their displacements
are at most $\ell+2$. Therefore
\[
 \clen{\alpha(v)}\le L(\ell+2)\le3L\ell
 =\frac{3L}{|p|}\clen{\alpha(u)}.
\]
This is specifically an ambient-automorphism argument. It completes
the one-sided criterion. Applying it in both directions gives
commensurability of the two oriented stabilizers.
